% 阶段报告五导读。编译：xelatex summary_plain_zh.tex（两遍）。
\documentclass[12pt,a4paper]{article}
\linespread{1.06}
\usepackage[margin=0.9in]{geometry}
\usepackage{amsmath,booktabs,tabularx,graphicx,xeCJK}
\setCJKmainfont{FandolSong-Regular}[Extension=.otf,BoldFont=FandolSong-Bold]
\setCJKsansfont{FandolHei-Regular}[Extension=.otf,BoldFont=FandolHei-Bold]
\usepackage[hidelinks]{hyperref}
\renewcommand{\figurename}{图}
\renewcommand{\tablename}{表}
\emergencystretch=3em
\setlength{\parskip}{3pt}
\title{\bfseries 分错的细胞核，是模型没学会，还是标签本身说不清？\\[7pt]
\large 阶段报告五导读}
\author{课程项目 bdd26}
\date{2026 年 10 月 13 日}
\begin{document}
\maketitle

\section*{先说结论}

上一阶段结束后，我们手里有一个奇怪的事实：模型找核的能力（用不区分类别的 bPQ 衡量，约 0.665）明显好于它认类别的能力（官方 mPQ 约 0.500）。这一轮我们沿着这条缝往下挖，得到两个结论：

\begin{itemize}
\item \textbf{“把已找到的核改对类别”这条路走不通。}如果请一位全知者把每个已匹配的核改成正确类别，mPQ 能从 0.4995 涨到 0.6632；但我们试了四种越来越复杂的自动改法，最好的一种只提高 0.0002，而且把最稀有的 Dead 类改得更差。
\item \textbf{标签本身有说不清的地方，而且有证据。}模型分错的核里，由 12 个“没见过这批图”的检测器组成的评审团有 \textbf{72.5\%} 站在模型的“错”答案一边，只有 18.2\% 支持标注答案；而模型分对的核上，评审团翻票率只有 1.9\%。最集中的模糊带在炎症（Inflammatory）与结缔（Connective）两类之间。
\end{itemize}

另外，我们原计划请人类专家盲审候选核，但专家资源不可得。改用“检测器评审团”后，得到一个附带发现：在标注完全没覆盖的位置，有 66--73\% 的候选核能获得多个不同检测器的一致检出——随机放置的假圆盘只有 0.75\% 的概率出现这种一致。

\section*{上限诊断：如果类别全对，能到多少？}

先做一个思想实验：保持核的数量和轮廓一动不动，只把每个与真值匹配上的核的类别改成真值类别，指标能到多少？这种用标准答案做的理想化试验叫 oracle。结果：三划分平均 mPQ 从 0.4995 到 \textbf{0.6632}；如果只改真值属于炎症或结缔的核，能到 0.6141，约占全部增益的七成。

这个数字是“几何现状下的类别天花板”，不是模型能达到的成绩。它只说明：类别方向值得查，而且炎症/结缔是第一嫌疑。（这个诊断读了历史上的测试预测，属于只读的审计操作；后面所有新实验都只用训练和验证数据。）

\section*{四种简单改法，全部无效}

我们在验证折上比较了四种不改几何的改类方法：现行多数票（T0）、实例内概率平均（T1）、把模型内部特征池化后接一个线性分类器（T2）、再加一圈周边上下文（T3）。事先写好的通过条件是 mPQ 至少提高 0.003、区间不含零、Dead 类不倒退超过 0.002 等。

结果是门没过。最好的 T2 只有 $+0.0002$，置信区间跨零；T3 是 $-0.0017$；T2 的 Dead 还下降了 0.0053，触发护栏。失败原因可以从三个准确率看出来：线性头在真值区域上判对 85.3\%，在预测实例上判对 85.8\%，而现行多数票本来就有约 85.7\%——\textbf{新头读出的信息和旧头用的一样多}。加上此前外部病理模型（CONCH）重分类也失败，事后改类这条路正式关闭。

\section*{请不到专家，怎么办？}

下一步本该请两名合格评审对 200 个窗口盲审（材料包已备好，评审规则事先冻结）。但专家资源确认不可得。经批准改为两件事：

\textbf{让检测器当评审（E1-C）。}13 个成员组成评审团：两种架构、多个随机种子、两个训练折，都不曾见过待审的验证折。每个候选核数一数有多少成员也检出了它（IoU$\geq$0.3 记一票）。为了给“多检测器一致”定个标准，我们做了个空对照：在真值核之外随机放 2,000 个同样大小的圆盘，它们平均只得到 0.07 票，得到 3 票以上的只有 0.75\%。

真实的候选核完全不同：不管 EF 规则接受还是拒绝，平均都有 5 票上下；与真值匹配的比例随票数单调上升（0--2 票约 9\%，5 票以上 45--55\%）。最有意思的是与标注\textbf{零重叠}的那批候选：其中 66--73\% 拿到 3 票以上、且票来自不同训练折和不同架构。多个独立训练的检测器在同一位置一致看到“像个核”的东西——这构成\textbf{标签可疑的证据}（可能是漏标，可能是标注单位不同，也可能大家共享某种误报模式；我们不下生物学结论）。

\textbf{让视觉大模型当评审（E1-M）。}两轮独立盲评互相之间的一致性极差（kappa 在 $-0.09$ 到 $+0.35$ 之间），其中一轮对几乎全为真核的对照组只判 40.9\% 为真。结论：模型评审目前不可用，这次只是验证了流程。

\section*{关键一问：评审团觉得“错”答案更像对的}

有了评审团，就可以问那个核心问题：模型\textbf{分错}的核，评审团怎么看？

把验证折上 6,858 个分错的已匹配核拿出来逐个投票（评审团去掉与被审模型同源的成员，剩 12 席）。规则事先冻结：得票最多的类需要总票数至少 3、严格领先、且票来自两个训练折。结果：\textbf{72.5\% 的核上，评审团的多数答案就是模型的“错”答案}；与标注一致的只有 18.2\%，两个比例的置信区间完全分开。作为对照，模型分对的 4 万多个核上，评审团改票的比率只有 1.9\%——说明评审团不是在机械复读某个模型。

分错误的类型看，模糊带非常集中：真值炎症、模型判结缔的 1,053 例中 74.8\% 支持模型的判断；反方向 78.5\%。唯一的例外是真值上皮、模型判肿瘤的格子，那里标注答案反而占优（51.6\% 对 30.1\%）。Dead 相关的格子方向相同但样本少，单独记录。

按事先登记的判读规则，这构成\textbf{“炎症/结缔标签边界模糊”的证据}。它同时解释了为什么简单改类失败：监督信号本身在这些边界上是含糊的，模型学到的就是数据教它的东西。这不是说 PanNuke 标注得差——炎症细胞和结缔细胞在 H\&E 染色下本来就看不清界限；也不是生物学结论——评审团与被审模型共享了部分设计。它只是说：\textbf{这部分缺口，靠在现有标签上继续训练是消不掉的。}

\section*{最后一试：让模型学会挑候选，还是不行}

E1-C 显示“票数”能预测候选核的好坏。于是我们训练了一个逻辑回归，用面积、形状、置信度等 12 个廉价特征（一路加上评审团票数）预测每个候选能否匹配真值。预测本身相当成功：交叉验证 AUC 0.721，加票数 0.737。

但把它接进实际流程后，按事先冻结的选择规则，三个学习臂全部退回“不加任何候选”。原因在一个具体的门槛上：选择规则要求 strict Dead（对“图里没有 Dead 却报出 Dead”也计罚的严格指标）不能低于不加候选的基线 0.12957，而所有学得的配置最高只有 0.12454。反过来说，学得的配置在官方口径上其实更好（有一个配置 Dead、bPQ、mPQ 三项全升），代价是 strict Dead 掉 0.013——\textbf{便宜的特征看得见“这个候选像不像核”，看不见“这张图里有没有这个类”}。按预注册规则，这条线停止，不加菜单、不换模型挽救。

失败留下一句有用的话：如果将来重开，该学的是\textbf{图像级}的缺类风险（比如这张图按组织类型看有没有 Dead），而不是更复杂的候选级打分。

\section*{哪些结果核对过，边界在哪里}

全部表格数字在写报告前从冻结产物重新导出并独立核验（核验脚本会重放选择规则、重算一致率，对不上就报错退出）；全量测试 335 项通过；两项实验的终版产物在独立评审后重跑，与第一版逐位一致。

边界也要讲清楚：所有新实验都在单一划分（split 1）的验证折上，基模型只有一个；“标签可疑”是统计证据，不能量化漏标率；评审团 12 席中 10 席与被审模型同属 CellViT 家族（两席 HoVer-Net 是仅有的跨架构对照，信号减弱但不反转）；本轮没有新的测试集成绩，也没有任何训练。人类专家评审的义务保留：材料包不变，将来若有两名合格评审，原协议可以直接执行。

下一步优先做的事，是把“标签模糊”写进最终报告的分析框架（组织$\times$类别的误差分解要和评审团证据对照呈现），以及评估图像级缺类风险特征是否值得立项。在这两件事之前，不启动新的训练。

完整数字、判定规则与复现命令见本目录的正式报告。
\end{document}
